\documentclass[11pt]{article}
\usepackage[margin=1in]{geometry}
\usepackage{amsmath,amssymb,amsthm}
\usepackage[T1]{fontenc}
\usepackage{lmodern}
\usepackage{graphicx}
\usepackage{hyperref}

\newtheorem{theorem}{Theorem}
\newtheorem{lemma}{Lemma}
\newtheorem{remark}{Remark}

\newcommand{\eps}{\varepsilon}
\newcommand{\R}{\mathbb R}
\newcommand{\cA}{\mathcal A}
\newcommand{\cL}{\mathcal L}
\newcommand{\perpB}{\perp_B}
\newcommand{\perpBe}{\perp_B^\varepsilon}
\newcommand{\dist}{\operatorname{dist}}

\title{A two-dimensional angular-locus classification for Property P}
\author{Codex agent lane 16}
\date{July 19, 2026}

\begin{document}
\maketitle

\section*{Source Question}

Mandal--Manna--Paul--Sain define \(x\perpBe y\) by
\[
  \|x+\lambda y\|^2\geq \|x\|^2-2\eps\|x\|\,\|\lambda y\|
  \qquad(\lambda\in\R)
\]
and say that a bounded linear operator \(T\) preserves \(\eps\)-orthogonality
at every point if
\[
  x\perpB y\quad\Longrightarrow\quad Tx\perpBe Ty
  \qquad(x,y\in X).
\]
Their Property P asks whether, for a given Banach space \(X\), some
\(\eps\in(0,1)\) forces every such \(T\in\cL(X)\) to be a scalar multiple of an
isometry \cite{MMPS}.  The source question and definition appear here:

\begin{center}
\includegraphics[width=0.88\linewidth]{figures/open_problem_crop.png}
\end{center}

This packet gives a claimed complete answer for real two-dimensional Banach
spaces.

\section*{Definitions and Tools}

For \(0\neq x\in X\), put
\[
  J(x)=\{f\in X^*: \|f\|=1,\ f(x)=\|x\|\}.
\]
The projective angular locus is
\[
  \cA_X=\{[x]\in\mathbb P(X): J(x)\text{ is not a singleton}\}.
\]
Thus \([x]\in\cA_X\) exactly when the norm is nonsmooth in the direction \(x\).

We use the source's support-functional characterization:
\[
  u\perpBe v
  \quad\Longleftrightarrow\quad
  \exists f\in J(u)\text{ such that } |f(v)|\leq\eps\|v\|
\]
\cite[Theorem 1.3]{MMPS}.  We also use the exact
Blanco--Koldobsky--Turnsek theorem: a linear operator preserving exact
Birkhoff--James orthogonality is a scalar multiple of an isometry
\cite{Koldobsky,BlancoTurnsek}.

Finally, we use Chmielinski--Stypka's bounded-below theorem for the same
approximate preservation relation.  If an additive map \(T:X\to Y\) between
real normed spaces satisfies
\[
  x\perpB y\Longrightarrow Tx\perpBe Ty,
\]
then \(T\) is linear and bounded and
\[
  (1-8\eps)\|T\|\,\|x\|\leq\|Tx\|\leq\|T\|\,\|x\|,
  \qquad x\in X
\]
\cite[Theorem 4.4]{CS}.  We apply only the displayed estimate, and only to
linear operators.

\section*{Proof Intuition}

The classification is governed by projective first-order rigidity.  Suppose
Property P fails.  Then there are norm-one nonsimilarities \(T_n\) preserving
\(\eps_n\)-orthogonality with \(\eps_n\downarrow0\).  The
Chmielinski--Stypka estimate makes a subsequential limit an isometry; composing
it away gives \(S_n\to I\).  Dividing \(S_n\) by its distance to the scalar
line gives a non-scalar first-order operator \(A\).

At a nonsmooth direction \(a\), the support set \(J(a)\) is a nontrivial
segment.  If \(A\) moves the projective line \([a]\) to first order, then
\([S_n a]\) approaches \([a]\) from one side.  Planar convexity says that all
support functionals at \(S_n a\) then collapse to the corresponding one-sided
endpoint of \(J(a)\).  Choosing an exact orthogonality direction from the
opposite endpoint gives a contradiction to approximate preservation.  Hence
every angular line must be an eigenline of \(A\).  Three projective eigenlines
force \(A\) to be scalar, contradicting the normalization.

The converse is also geometric.  If there are at most two angular lines, a
small diagonal perturbation can fix all of them while stretching one transverse
direction.  Such perturbations are not similarities, but because they fix the
only possible support discontinuities, they preserve exact orthogonality up to
an arbitrarily small Chmielinski error.

\begin{lemma}[One-sided support collapse in a normed plane]
\label{lem:one-sided}
Let \(X\) be a real two-dimensional Banach space, let \(a\in S_X\), and suppose
\(J(a)\) is not a singleton.  Choose one of the two local sides of the
projective line \([a]\).  There is an endpoint \(j_+\) of the segment \(J(a)\)
such that whenever \(x_n\in S_X\) tends projectively to \(a\) from that side
and \(f_n\in J(x_n)\), every norm-convergent subsequence of \(f_n\) converges
to \(j_+\).  The other side gives the other endpoint \(j_-\).
\end{lemma}

\begin{proof}
This is the standard one-dimensional subdifferential fact for a planar convex
body.  We include the reduction for reference.  In an affine chart around
\(a\), the boundary of \(B_X\) on either side of \(a\) can be represented as
the graph of a concave function \(\phi\) of one real variable.  Supporting
functionals at a boundary point correspond exactly to slopes in the
subdifferential interval of \(\phi\) at the corresponding parameter.  For a
concave function on an interval, the left and right derivative functions are
monotone, and the subdifferential intervals \(\partial\phi(t_n)\) converge to
the appropriate one-sided derivative as \(t_n\to0^+\) or \(t_n\to0^-\).  Under
the affine identification this one-sided derivative is precisely one endpoint
of \(J(a)\).  This proves the claim.
\end{proof}

\begin{lemma}[Infinitesimal preservers fix angular lines]
\label{lem:infinitesimal-fixes}
Let \(X\) be a real two-dimensional Banach space.  Suppose
\(\eps_n\downarrow0\), \(S_n\to I\), and each \(S_n\) preserves
\(\eps_n\)-orthogonality at every point.  Assume that
\[
  S_n=\lambda_n I+d_n A_n,\qquad \lambda_n\to1,\quad d_n\downarrow0,\quad
  A_n\to A.
\]
If \([a]\in\cA_X\), then \(Aa\in\R a\).
\end{lemma}

\begin{proof}
Suppose \(Aa\notin\R a\) for some unit representative \(a\).  After passing to
a subsequence, the projective points \([S_n a]\) approach \([a]\) from one
fixed side.  Let \(j_+\) be the corresponding endpoint of \(J(a)\) from
Lemma~\ref{lem:one-sided}.  Choose another support functional
\(j_0\in J(a)\), \(j_0\neq j_+\), and choose \(0\neq y\in\ker j_0\).  Then
\(a\perpB y\), while \(j_+(y)\neq0\), because two distinct support functionals
at \(a\) have distinct kernels in a plane.

Since \(S_n\) preserves \(\eps_n\)-orthogonality, there is
\(g_n\in J(S_n a)\) such that
\[
  |g_n(S_n y)|\leq \eps_n\|S_n y\|.
\]
Passing to a subsequence, \(g_n\) converges in \(S_{X^*}\).  By
Lemma~\ref{lem:one-sided}, the only possible limit is \(j_+\).  Letting
\(n\to\infty\) in the displayed inequality gives \(j_+(y)=0\), a
contradiction.  Therefore \(Aa\in\R a\).
\end{proof}

\begin{lemma}[Three angular lines imply Property P]
\label{lem:three-positive}
If a real two-dimensional Banach space \(X\) has at least three projective
angular lines, then \(X\) has Property P.
\end{lemma}

\begin{proof}
Assume Property P fails.  Then for some \(\eps_n\downarrow0\) there are
norm-one operators \(T_n\) preserving \(\eps_n\)-orthogonality at every point
which are not scalar multiples of isometries.  By the Chmielinski--Stypka
estimate,
\[
  (1-8\eps_n)\|x\|\leq\|T_nx\|\leq\|x\|.
\]
After passing to a subsequence, \(T_n\to U\), and the displayed estimate shows
that \(U\) is a linear isometry.  Replacing \(T_n\) by \(S_n=U^{-1}T_n\), we
may assume \(S_n\to I\).  Isometric conjugation preserves the same approximate
orthogonality condition, and \(S_n\) is still not a scalar multiple of an
isometry.

Let
\[
  d_n=\dist(S_n,\R I)>0
\]
and choose \(\lambda_n\in\R\) with \(\|S_n-\lambda_n I\|=d_n\).  Since
\(S_n\to I\), we have \(\lambda_n\to1\) and \(d_n\to0\).  Put
\[
  A_n=d_n^{-1}(S_n-\lambda_n I).
\]
After a subsequence, \(A_n\to A\).  Moreover \(\dist(A,\R I)=1\), so \(A\) is
not scalar.

Lemma~\ref{lem:infinitesimal-fixes} says that every angular projective line is
an eigenline of \(A\).  But a non-scalar linear operator on a two-dimensional
real vector space has at most two projective eigenlines.  This contradicts the
existence of three angular lines.  Hence Property P holds.
\end{proof}

\begin{lemma}[Finite angular stabilizer perturbations]
\label{lem:finite-stabilizer}
Let \(X\) be a real two-dimensional Banach space with finite angular locus
\(\cA_X\).  Let \(R_n\to I\) be invertible operators such that every line in
\(\cA_X\) is invariant under every \(R_n\).  Then for every \(\eps>0\), all
sufficiently large \(n\) have the property that \(R_n\) preserves
\(\eps\)-orthogonality at every point.
\end{lemma}

\begin{proof}
It is enough to prove that the optimal Chmielinski error tends to zero.  If
not, there are \(\eta>0\), unit vectors \(x_n,y_n\), and a subsequence, still
called \(R_n\), such that \(x_n\perpB y_n\) but
\[
  \min_{g\in J(R_nx_n)}\frac{|g(R_ny_n)|}{\|R_ny_n\|}\geq\eta.
\]
Pass to a further subsequence with \(x_n\to x\) and \(y_n\to y\).  The exact
orthogonality relation is closed, so \(x\perpB y\).

For each \(n\), choose \(f_n\in J(x_n)\) with \(f_n(y_n)=0\).  We will find
\(g_n\in J(R_nx_n)\) with \(g_n(R_ny_n)\to0\), contradicting the displayed
lower bound.

If \(x\) is smooth, then the support-functional map is continuous in a
neighborhood of \(x\) avoiding the finite set \(\cA_X\), so the unique
functionals in \(J(R_nx_n)\) converge to the same limit as \(f_n\).  Hence
\(g_n(R_ny_n)\to0\).

Now suppose \([x]\in\cA_X\).  Because \(\cA_X\) is finite, a punctured
projective neighborhood of \([x]\) contains no other angular line.  After
passing to a subsequence, either \(x_n\) lies on \([x]\) for all \(n\), or
\([x_n]\) approaches \([x]\) from one fixed side.  In the first case,
\([R_nx_n]=[x]\), and for large \(n\) the scalar on this invariant line is
positive; taking \(g_n=f_n\) gives \(g_n\in J(R_nx_n)\) and
\[
  g_n(R_ny_n)=g_n(y_n)+g_n((R_n-I)y_n)\to0.
\]
In the second case, \(R_n\) fixes \([x]\) and is close to the identity, so
\([R_nx_n]\) approaches \([x]\) from the same side.  By the one-sided support
collapse lemma, both \(f_n\) and the unique member \(g_n\in J(R_nx_n)\) tend to
the same endpoint of \(J(x)\).  Therefore
\[
  |g_n(R_ny_n)|\leq |(g_n-f_n)(R_ny_n)|+|f_n((R_n-I)y_n)|+|f_n(y_n)|\to0.
\]
This final contradiction proves the lemma.
\end{proof}

\begin{lemma}[At most two angular lines imply failure of Property P]
\label{lem:two-negative}
If a real two-dimensional Banach space \(X\) has at most two projective angular
lines, then \(X\) does not have Property P.
\end{lemma}

\begin{proof}
Let \(L_1,L_2\) be two independent projective lines containing \(\cA_X\); if
\(\cA_X\) has fewer than two elements, choose the missing line or lines
arbitrarily.  Pick unit vectors \(e_i\in L_i\).  For \(t>0\), define
\[
  R_t e_1=e_1,\qquad R_t e_2=(1+t)e_2.
\]
Then \(R_t\to I\), and \(R_t\) leaves every angular line invariant.  By
Lemma~\ref{lem:finite-stabilizer}, for every \(0<\eps<1\), all sufficiently
small \(t>0\) make \(R_t\) preserve \(\eps\)-orthogonality at every point.

However \(R_t\) is not a scalar multiple of an isometry for \(t\neq0\).  Indeed,
if \(R_t=cU\) with \(U\) an isometry, then from the unit vector \(e_1\) we get
\(|c|=1\), while from \(e_2\) we get \(|c|=1+t\), impossible.  Thus no
\(\eps\) can witness Property P.
\end{proof}

\begin{theorem}[Two-dimensional classification]
\label{thm:classification}
Let \(X\) be a real two-dimensional Banach space.  Then \(X\) has Property P if
and only if its projective angular locus \(\cA_X\) contains at least three
projective lines.
\end{theorem}

\begin{proof}
The positive implication is Lemma~\ref{lem:three-positive}.  The negative
implication is Lemma~\ref{lem:two-negative}.
\end{proof}

\section*{Consequences}

For a two-dimensional polyhedral space, the angular lines are exactly the
projective vertex lines.  Hence the theorem recovers the previous polyhedral
classification: a polygonal unit ball has Property P exactly when it has at
least six extreme points.  The theorem also handles non-polyhedral examples
with infinitely many corners; any such plane has Property P as soon as it has
three distinct angular directions.

Smooth two-dimensional spaces fail Property P because \(\cA_X=\varnothing\).
This agrees with the separate lane-16 uniformly smooth packet.

\section*{Verification and Novelty Check}

Local indexes were checked for arXiv:2501.10719, Property P, two-dimensional
polyhedral classifications, finite angular locus, and Birkhoff--James
orthogonality.  Web searches on July 19, 2026 located the source arXiv record,
the published Colloquium Mathematicum record, the 2026 Lebesgue--Bochner
continuation, and Chmielinski--Stypka's 2025 bounded-below theorem.  No exact
two-dimensional angular-locus classification was found.

Human-review recommendation: focus first on Lemma~\ref{lem:one-sided} and its
use in Lemmas~\ref{lem:infinitesimal-fixes} and \ref{lem:finite-stabilizer}.
Those are the only genuinely delicate planar convexity steps.  The remaining
argument is compactness, first-order projective rigidity, and elementary
diagonal perturbations.

\begin{thebibliography}{9}

\bibitem{MMPS}
K. Mandal, J. Manna, K. Paul, and D. Sain,
\emph{On approximate preservation of orthogonality and its application to
isometries}, arXiv:2501.10719; Colloquium Mathematicum 179 (2025), 189--211.
DOI: \href{https://doi.org/10.4064/cm9612-7-2025}{10.4064/cm9612-7-2025}.

\bibitem{CS}
J. Chmielinski and R. Stypka,
\emph{Additive operators approximately preserving Birkhoff--James
orthogonality}, Aequationes Mathematicae 99 (2025), 2847--2854.
DOI: \href{https://doi.org/10.1007/s00010-025-01210-4}{10.1007/s00010-025-01210-4}.

\bibitem{Koldobsky}
A. Koldobsky,
\emph{Operators preserving orthogonality are isometries}, Proc. Roy. Soc.
Edinburgh Sect. A 123 (1993), 835--837.

\bibitem{BlancoTurnsek}
A. Blanco and A. Turnsek,
\emph{On maps that preserve orthogonality in normed spaces}, Proc. Roy. Soc.
Edinburgh Sect. A 136 (2006), 709--716.

\bibitem{Rockafellar}
R. T. Rockafellar,
\emph{Convex Analysis}, Princeton University Press, 1970.

\bibitem{MJ}
Mohit and R. Jain,
\emph{Approximate Birkhoff--James orthogonality preserver on
Lebesgue--Bochner spaces}, arXiv:2601.02786, 2026.

\end{thebibliography}

\end{document}
